% ========================================================================
% فصل ۴: حملات امنیتی بر روی شبکه اقتضایی خودرویی
% ========================================================================
\section{حملات امنیتی بر روی شبکه اقتضایی خودرویی}

\subsection{الگوی تهدید}

ماهیت پخشی ارتباطات بی‌سیم، شبکه اقتضایی خودرویی را در برابر طیف گسترده‌ای از حملات آسیب‌پذیر می‌کند. از آنجا که اطلاعات میان اشیاء \lr{V2X} به‌صورت پخشی رد و بدل می‌شود، مهاجم می‌تواند آن‌ها را شنود، جعل یا بازپخش کند و حتی خود را به جای یک موجودیت مجاز، مانند یک خودرو یا واحد کنارجاده‌ای، جا بزند و داده نادرست ارسال کند \cite{farsimadanReviewSecurityChallenges2025}. پیامد این حملات می‌تواند بسیار جدی باشد؛ مهاجم ممکن است با سوءاستفاده از آسیب‌پذیری‌های سامانه به خودروها دسترسی یابد و آن‌ها را کنترل کند و شرایط رانندگی خطرناک و تصادفات مرگبار به وجود آورد \cite{farsimadanReviewSecurityChallenges2025}. در همین حال، تحلیل تهدیدات و معماری امنیتی این شبکه‌ها به دنبال حفظ تعادل میان کاهش مؤثر تهدیدات و حداقل تأثیر بر عملکرد سامانه و راحتی کاربر است \cite{SecuringVehicularAd,VanetSecurityPrivacy}.

برای تحلیل نظام‌مند این تهدیدها، \lr{Raya} و \lr{Hubaux} مدل مهاجمی با چهار بعد پیشنهاد کرده‌اند. در بعد نخست، مهاجم داخلی عضو احرازهوید‌شدهٔ شبکه است و کلید عمومی تأییدشده دارد، در حالی که مهاجم خارجی از دید اعضا یک مزاحم است و دامنهٔ حملاتش محدودتر است. در بعد دوم، مهاجم بدخواه بدون توجه به هزینه و بدون نفع شخصی آسیب می‌زند، اما مهاجم عقلانی به دنبال سود شخصی است و از این رو رفتارش قابل پیش‌بینی‌تر است. در بعد سوم، مهاجم فعال بسته یا سیگنال تولید می‌کند و مهاجم غیرفعال تنها کانال بی‌سیم را شنود می‌کند. در بعد چهارم، مهاجم محلی دامنهٔ محدودی دارد، اما مهاجم گسترده موجودیت‌هایی را در سراسر شبکه کنترل می‌کند \cite{SecuringVehicularAd}. این نویسندگان فرض می‌کنند که به دلیل تعداد زیاد اعضای مستقل و حضور عامل انسانی، رفتار نادرست به احتمال زیاد رخ خواهد داد؛ بنابراین سطح اعتماد به خودروها و حتی به ایستگاه‌های پایهٔ ارائه‌دهندگان خدمات پایین در نظر گرفته می‌شود \cite{SecuringVehicularAd}.

در ادامه، حملات بر اساس هدف و لایهٔ مورد هجوم طبقه‌بندی می‌شوند. این طبقه‌بندی با طبقه‌بندی پنج‌گانهٔ مرور فارسیمدان و همکاران — حملات بر زیرساخت، بر نرم‌افزار و سخت‌افزار، بر حریم خصوصی، بر اعتماد داده و مبتنی بر الگوی رفتاری \cite{farsimadanReviewSecurityChallenges2025} — و با حملات پنج‌گانهٔ مرور سیلا و همکاران — شنود، جعل، حملهٔ منع سرویس، سیبل و فردی در میان \cite{VanetSecurityPrivacy} — هماهنگ است.

\subsection{حملات بر دسترس‌پذیری}

الزام دسترس‌پذیری موجب می‌شود ارتباطات میان اشیاء \lr{V2X} تحلیل و به‌موقع در دسترس گیرندگان مورد نظر قرار گیرند، و این الزام استفاده از تکنیک‌های رمزنگاری سبک و کم‌هزینه را ضروری می‌کند \cite{farsimadanReviewSecurityChallenges2025}. نخستین گروه حملات، حملات منع سرویس \LTRfootnote{\lr{Denial of Service (DoS)}} و نسخهٔ توزیع‌شدهٔ آن است. در این حملات، یک هماهنگ‌کننده که تعداد زیادی ربات کنترل می‌کند، شبکه را با بسته‌های دشمن‌واره یا حداقل ناخواسته اشباع می‌کند \cite{farsimadanReviewSecurityChallenges2025} و ارتباطات مشروع را مانع می‌شود؛ نمونه‌ای از آن سیل کانال کنترل است \cite{farsimadanReviewSecurityChallenges2025}. مرور سیلا و همکاران می‌افزاید که چنین حملاتی با سیلی از درخواست‌های غیرضروری عملکرد شبکه را به سختی متحمل می‌شوند یا کانال‌های ارتباطی را کاملاً از کار می‌اندازند و می‌توانند به وضعیت‌های خطرناک در جاده منجر شوند \cite{VanetSecurityPrivacy}. در برابر این حملات، تصدیق قوی و امضای دیجیتال در سطح گره \cite{farsimadanReviewSecurityChallenges2025}، محدودسازی نرخ، تشخیص ناهنجاری و بخش‌بندی شبکه \cite{VanetSecurityPrivacy} پیشنهاد می‌شوند.

پارازیت \LTRfootnote{\lr{Jamming}} گونهٔ دیگری از حمله به دسترس‌پذیری است که در لایهٔ فیزیکی کانال‌های بی‌سیم را مختل می‌کند؛ مرور فارسیمدان و همکاران مدلی برای تشخیص پارازیت فرکانس رادیویی در ارتباطات خودرو با همه‌چیز بدون نظارت گزارش می‌کنند \cite{farsimadanReviewSecurityChallenges2025}. سیل پیام \LTRfootnote{\lr{Flooding}} نیز با طوفان پخش پیام، پهنای باند را اشباع می‌کند و از جمله حملاتی است که مینتمور و سن آن را در مجاورت گودال سیاه، رهاکردن بسته و اطلاعات نادرست بر روی پروتکل‌های مسیریابی \lr{AODV} و \lr{GPSR} شبیه‌سازی کرده‌اند \cite{VanetSecurityPrivacy}.

\subsection{حملات بر محرمانگی}

استراق سمع \LTRfootnote{\lr{Eavesdropping}} حمله‌ای غیرفعال است که در آن ارتباط میان خودروها به‌صورت غیرمجاز رهگیری می‌شود و محرمانگی و حریم خصوصی داده‌های منتقل‌شده را به خطر می‌اندازد \cite{VanetSecurityPrivacy}. داده‌هایی که در معرض خطر قرار می‌گیرند شامل جزئیات موقعیت، الگوهای سفر و داده‌های شخصی سرنشینان است \cite{VanetSecurityPrivacy}. در طبقه‌بندی رفتاری مرور فارسیمدان و همکاران، شنود به دستهٔ رفتار خودخواهانه تعلق دارد و به‌عنوان یک حملهٔ منفعل شناخته می‌شود \cite{farsimadanReviewSecurityChallenges2025}. این حمله در پروتکل‌های ورود بدون کلید از راه دور، ارتباطات فوق‌پهن‌باند با سکوهای کم‌توان و شبکه‌های سلولی نیز گزارش شده است \cite{farsimadanReviewSecurityChallenges2025}. در برابر آن، رمزنگاری قوی، پروتکل‌های تبادل کلید امن و بخش‌بندی شبکه برای محدودکردن دامنهٔ جغرافیایی هر ارتباط مؤثرند \cite{VanetSecurityPrivacy}.

تحلیل ترافیک \LTRfootnote{\lr{Traffic analysis}} گام فراتری است و از روی ترافیک مشاهده‌شده الگوهای رفتاری و مکانی استنباط می‌کند. این حمله در مرور فارسیمدان و همکاران جزو رفتار خودخواهانه و منفعل دسته‌بندی شده و دفاع‌هایی چون کلیدهای ناشناس پویا، شبه‌نام و امضای گروهی برای آن برشمرده می‌شود \cite{farsimadanReviewSecurityChallenges2025}. توجه به این نکته ضروری است که هیچ‌یک از سه منبع مرورشده در این گزارش حملات کانال جانبی \LTRfootnote{\lr{Side-channel}} — یعنی استخراج کلید از راه نشت زمان‌بندی، توان مصرفی یا تابش الکترومغناطیسی — را برای این شبکه‌ها پوشش نداده‌اند و از این رو آن‌ها در اینجا توصیف نمی‌شوند.

\subsection{حملات بر صحت، یکپارچگی و تصدیق هویت}

جعل پیام \LTRfootnote{\lr{Spoofing/Forgery}} و جعل هویت \LTRfootnote{\lr{Impersonation/Masquerade}} از نخستین حملات علیه صحت و اصالت هستند. در جعل پیام، مهاجم اطلاعات نادرست یا جعلی تولید می‌کند \cite{farsimadanReviewSecurityChallenges2025}، مثلاً اطلاعات موقعیت نادرست \cite{farsimadanReviewSecurityChallenges2025}. مرور سیلا و همکاران می‌افزاید که مهاجمان خود را به جای خودرو یا جزء زیرساختی جا می‌زنند تا اطلاعات یا دستور نادرست ارسال کنند و این امر می‌تواند آشوب در سامانهٔ مدیریت ترافیک ایجاد کند؛ پیامدهای آن شامل اطلاعات مسیریابی نادرست، هشدارهای نادرست سامانهٔ اورژانس یا به‌روزرسانی‌های جعلی ترافیک است \cite{VanetSecurityPrivacy}. در طبقهٔ حملات زیرساختی، جعل هویت می‌تواند به ساختن گودال سیاه یا تقلید خودروی امدادی منجر شود و دفاع رایج آن عدم انکار است \cite{farsimadanReviewSecurityChallenges2025}. در برابر جعل پیام، امضای پیام‌های هشدار، زیرساخت کلید عمومی خودرویی، ارتباط گروهی، بررسی پلاسیبیلتی و رادار درون‌خودرو یا گواهی‌های رمزنگاری پیشنهاد می‌شود \cite{farsimadanReviewSecurityChallenges2025}.

دستکاری پیام \LTRfootnote{\lr{Message tampering}} شامل تغییر، حذف، اصلاح و ساختن داده‌های از پیش موجود است \cite{farsimadanReviewSecurityChallenges2025} و در برابر آن همبستگی داده و چالش و پاسخ با طرح امضای گروهی مؤثر است \cite{farsimadanReviewSecurityChallenges2025}. حملهٔ پخش مجدد \LTRfootnote{\lr{Replay}} از پیام‌های معتبر ضبط‌شده به‌صورت نادرست استفاده می‌کند و مقابله با آن با برچسب زمانی و شمارهٔ ترتیب امضاشده \cite{farsimadanReviewSecurityChallenges2025} یا با رمزنگاری متقارن با افشای تأخیری کلیدها (\lr{TESLA++}) انجام می‌شود \cite{farsimadanReviewSecurityChallenges2025}.

حملهٔ فردی در میان \LTRfootnote{\lr{Man-in-the-Middle (MiM)}} مهاجمی را میان فرستنده و گیرنده جا می‌زند و «اصالت، عدم انکار و یکپارچگی سامانه‌های خودرویی را به‌شدت تحت تأثیر قرار می‌دهد» \cite{farsimadanReviewSecurityChallenges2025}. در این حمله، مهاجم ارتباط میان دو طرف را بدون اطلاع آن‌ها رهگیری و احتمالاً تغییر می‌دهد و می‌تواند گیرندگان را فریب دهد یا رفتار خودروها را تغییر دهد \cite{VanetSecurityPrivacy}. دفاع‌ها شامل تصدیق قوی \cite{farsimadanReviewSecurityChallenges2025}، تصدیق همکارانهٔ دسته‌ای امضاها \cite{farsimadanReviewSecurityChallenges2025}، ترکیب رمزنگاری متقارن و نامتقارن برای یکپارچگی و محرمانگی \cite{VanetSecurityPrivacy} و پایش شبکه و تشخیص ناهنجاری است؛ مرور سیلا و همکاران همچنین استدلال می‌کنند که ماهیت غیرمتمرکز بلاکچین با حذف نیاز به مرجع مرکزی، خطر حملهٔ فردی در میان را کاهش می‌دهد و داده‌ها را در دفتری امن و تغییرناپذیر ثبت می‌کند \cite{VanetSecurityPrivacy}.

جعل موقعیت و دادهٔ حسگر گروهی از حملات اعتماد داده‌اند که ادراک نادرست ایجاد می‌کنند. در حملهٔ خودروی پنهان \LTRfootnote{\lr{Hidden vehicle}}، هشدار موقعیت جعلی یا جعل \lr{GPS} رخ می‌دهد و مدیریت اعتماد دفاع آن است \cite{farsimadanReviewSecurityChallenges2025}. در حملهٔ توهم \LTRfootnote{\lr{Illusion attack}}، اطلاعات نادرست با یک اتصال مشروع تولید می‌شود و تشخیص آن دشوار است؛ سامانهٔ موقعیت با امضا و ردیابی تفاضلی برای آن پیشنهاد شده‌اند \cite{farsimadanReviewSecurityChallenges2025}. حملهٔ ربایش نشست \LTRfootnote{\lr{Session hijacking}} نیز از گروه حملات زیرساختی است: تصدیق در آغاز نشست انجام می‌شود و مهاجم بعداً نشست را تصاحب می‌کند؛ دفاع، زیرساخت کلید عمومی و مرجع قابل اعتماد است که در آن واحد کنارجاده‌ای هویت خودرو را نزد مرجع بررسی و سپس کلید را با آن به اشتراک می‌گذارد \cite{farsimadanReviewSecurityChallenges2025}. سرانجام، کد موذیانه \LTRfootnote{\lr{Malicious code}} از حملات رفتار بدخواهانه است \cite{farsimadanReviewSecurityChallenges2025} و حملهٔ جست‌وجوی استردادی \LTRfootnote{\lr{Brute force}} به‌دلیل اتصال کوتاه‌مدت دشوار اما ممکن است و دفاع آن کلید و رمزنگاری قوی است \cite{farsimadanReviewSecurityChallenges2025}.

\subsection{حملات بر حریم خصوصی}

حملات بر حریم خصوصی مستقیماً به الزامات گمنامی، عدم ردیابی‌پذیری و پیوندناپذیری هدف می‌گیرند. در طبقه‌بندی مرور فارسیمدان و همکاران، افشای هویت \LTRfootnote{\lr{Identity disclosure}} نقض اطلاعات هویتی سرنشینان خودرو است و دفاع آن احراز هویت حافظ حریم خصوصی است \cite{farsimadanReviewSecurityChallenges2025}. ردیابی موقعیت و مسیر \LTRfootnote{\lr{Location tracking}} موقعیت خودرو و مسیری را که در بازهٔ زمانی طی کرده دنبال می‌کند و دفاع آن کلیدهای موقت و ناشناس است \cite{farsimadanReviewSecurityChallenges2025}.

حملهٔ پیوندپذیری \LTRfootnote{\lr{Linkability attack}} چند شبه‌نام یا پیام را به یک خودروی واحد متصل می‌کند. مرور ره‌آواتی آگوستینا و همکاران هشدار می‌دهد که راهبرد تغییر شبه‌نام با زمان ثابت که چگالی ترافیک را نادیده بگیرد، به حملهٔ پیوند می‌انجامد \cite{rahmawatiagustinaSystematicLiteratureReview2025}. نمایه‌سازی \LTRfootnote{\lr{Profiling}} امکان ساختن نمایهٔ رفتاری و هویتی در طول زمان را فراهم می‌کند؛ هم‌رسانی مداوم داده در شبکه به مهاجم اجازه می‌دهد از خودروها نمایه‌سازی کند و بدین‌سان حریم خصوصی مالکان آن‌ها را به خطر بیندازد \cite{rahmawatiagustinaSystematicLiteratureReview2025} و حریم خصوصی باید داده‌های کاربر را در برابر ردیابی و نمایه‌سازی محافظت کند \cite{VanetSecurityPrivacy}. دسترسی غیرمجاز نیز کنترل دسترسی را دور می‌زند؛ در برابر آن، شبه‌نام دفاع رایج است، به‌طوری‌که هر گرهٔ خودرو چند جفت کلید دارد و نمی‌توان گره را به شبه‌نام‌ها پیوند زد، اما مرجع مجاز به آن‌ها دسترسی دارد و خودروها پیش از انقضای شبه‌نام قبلی از طریق واحدهای کنارجاده‌ای شبه‌نام به‌روز دریافت می‌کنند \cite{farsimadanReviewSecurityChallenges2025}.

\subsection{حملات رفتاری گره‌های داخلی}

گره‌های داخلی که گواهی معتبر دارند، می‌توانند رفتارهای موذی‌ای از خود نشان دهند. رفتار خودخواهانه \LTRfootnote{\lr{Selfish behavior}} شامل جعل پیام، تحلیل ترافیک (منفعل)، شنود (منفعل) و انکار است و دفاع در برابر انکار، سخت‌افزار مورداعتماد است \cite{farsimadanReviewSecurityChallenges2025}. انکار \LTRfootnote{\lr{Repudiation}} یعنی انکار ارسال یا دریافت پیام‌ها؛ الزام عدم انکار اثبات ارتباط و تراکنش‌های داده را برای جلوگیری از انکار مشارکت طرفین فراهم می‌کند و ابزارهای آن شامل امضای دیجیتال، \lr{PKI}، مهر زمانی و ثبت وقایع امن است که کاربرد حقوقی پس از تصادف دارد \cite{VanetSecurityPrivacy}.

تبانی \LTRfootnote{\lr{Collusion}} تقلب هماهنگ میان چند گره برای فریب سیستم است. رایا و هوبو مفهوم مهاجمانی را با تعداد \lr{m} و \lr{n} گره تحت کنترل و پوشش تبانی در مدل خود گنجانده‌اند و نمونه‌ای از حملهٔ اطلاعات نادرست با تبانی را ترسیم می‌کنند \cite{SecuringVehicularAd}. در سامانهٔ مدیریت اعتماد مقیاس‌پذیر مبتنی بر زنجیره بلوکی نیز واحدهای کنارجاده‌ای برای پیشگیری از تبانی، هویت گذرای خودتولید به‌جای هویت دائمی یا زیرساخت کلید عمومی به کار می‌برند \cite{ghovanlooyghajarSBTMSScalableBlockchain2021}.

در حملهٔ گودال سیاه \LTRfootnote{\lr{Blackhole}}، بسته‌های عبوری رها می‌شوند و این رفتار از گونهٔ بدخواهانه است و دفاع آن تشخیص نفوذ ترکیبی و مسیریابی امن است \cite{farsimadanReviewSecurityChallenges2025}. حملهٔ سیبل \LTRfootnote{\lr{Sybil attack}} نیز در آن یک گره به‌غیرمجاز چندین هویت برای خود می‌سازد و می‌تواند سامانه‌های اعتماد و شهرت را با جابه‌جایی تصمیمات اجماعی یا اکثریتی به‌شخت خلل بزند \cite{VanetSecurityPrivacy}؛ در طبقهٔ رفتار بدخواهانه، سیبل «خودروهای متعددی با شناسه‌های مشابه در جاده می‌سازد» و دفاع آن بررسی موقعیت منطقی، گواهی، طول عمر پیام و گزارش به مرجع گواهی‌دهی است \cite{farsimadanReviewSecurityChallenges2025}. پیامدهای آن اختلال در سازوکارهای اجماع، سامانه‌های شهرت و پروتکل‌های مسیریابی و در نتیجه گزارش‌های جعلی ترافیک، جریان ترافیک دستکاری‌شده و حتی انزوای خودروهای مشروع از شبکه است \cite{VanetSecurityPrivacy}. از سوی دیگر، نبود مرحلهٔ لغو در چرخهٔ حیات شبه‌نام خطر حملهٔ سیبل و دسترسی غیرمجاز را بالا می‌برد \cite{rahmawatiagustinaSystematicLiteratureReview2025}.

\subsection{حملات لایه مسیریابی و شبکه}

حملات لایهٔ مسیریابی توانایی هدایت داده را در شبکهٔ چندپرشی هدف می‌گیرند. مینتمور و سن (۲۰۱۷) چهار حملهٔ گودال سیاه، رهاکردن بسته، سیل و اطلاعات نادرست را بر روی پروتکل‌های مسیریابی \lr{AODV} و \lr{GPSR} شبیه‌سازی کرده‌اند \cite{VanetSecurityPrivacy}. این حملات با نتایج حملات رفتاری گره‌های داخلی هم‌پوشانی دارند؛ به‌عنوان مثال، حملهٔ سیبل می‌تواند پروتکل‌های مسیریابی را نیز مختل کند \cite{VanetSecurityPrivacy} و مسیریابی داده یکی از سه مسئلهٔ بنیادین ارتباط میان‌خودرویی به شمار می‌رود \cite{hozouriOverviewVANETVehicular2023}. تکنیک‌های مسیریابی مبتنی بر توپولوژی و مبتنی بر موقعیت نیز در اینترنت خودروها بررسی شده‌اند \cite{xuInternetVehiclesBig2018}.

\subsection{حملات سخت‌افزاری، فیزیکی و فناوری‌های ارتباطی}

دستکاری سخت‌افزار \LTRfootnote{\lr{Hardware tampering}} توسط افراد مخرب، مثلاً هنگام سرویس سالانهٔ خودرو، صورت می‌گیرد و دفاع آن ماژول سکوی مورداعتماد است \cite{farsimadanReviewSecurityChallenges2025}. حملات ورود بدون کلید از راه دور \LTRfootnote{\lr{RKE}} شامل شنود، پارازیت، فردی در میان، جست‌وجوی استردادی و پخش مجدد است و آسیب‌پذیری \lr{Hitag2} نیز گزارش شده است؛ رمز ناشناختهٔ مخفی، گیرندهٔ مقاوم در برابر دستکاری و گذار به \lr{AES} دفاع‌های پیشنهادی هستند \cite{farsimadanReviewSecurityChallenges2025}. در حملهٔ دوقلوی شیطانی \LTRfootnote{\lr{Evil Twin}}، نقطهٔ دسترسی \lr{Wi-Fi} جعل می‌شود تا خودرو را فریب دهد؛ نمونهٔ هک تسلا مدل اس با رمز \lr{SSID} به‌صورت متن ساده، پارازیت، حملهٔ منع سرویس بر \lr{WPA} و دفاع با ترکیب تصادفی رمزشده و دروازهٔ جداگانه برای شناسایی دوقلوی شیطانی در همین حوزه گزارش شده‌اند \cite{farsimadanReviewSecurityChallenges2025}.

فناوری‌های ارتباطی دیگر نیز تهدیدهای ویژه‌ای دارند. در \lr{ZigBee}، افشای پیکربندی، سامانه‌های رمزنشده و پخش مجدد گزارش شده و دفاع‌ها شامل تشخیص نفوذ، کلید شبکهٔ پیش‌نصب، مهر زمانی و چرخش شناسهٔ دستگاه است \cite{farsimadanReviewSecurityChallenges2025}. در ارتباطات سلولی، نبود حفاظت رمزنگاری سیگنالینگ \lr{LTE}، شنود، پارازیت و ایستگاه‌های پایهٔ فمتوسل یا بی‌مجوز از تهدیدات هستند و \lr{SDN}/\lr{NFV} و مدیریت کلید سبک دفاع‌های پیشنهادی‌اند \cite{farsimadanReviewSecurityChallenges2025}. در ارتباطات فوق‌پهن‌باند نیز شنود در سکوهای کم‌توان گزارش شده و به‌جای رمزنگاری، پرش زمانی پیشنهاد می‌شود \cite{farsimadanReviewSecurityChallenges2025}.

\subsection{حملات زیرساخت و اعتماد}

زیرساخت مرکزی، چنان‌که فصل‌های پیش نشان داد، نقطهٔ شکست منفرد \LTRfootnote{\lr{Single Point of Failure (SPoF)}} است. در مدل استاندارد، عدم در دسترس بودن مرجع قابل اعتماد بر همهٔ موجودیت‌ها اثر می‌گذارد و چون این مرجع هویت همهٔ نهادها را می‌داند، هدف اصلی نقض حریم خصوصی می‌شود و تسخیر آن می‌تواند نقض گستردهٔ حریم خصوصی به همراه داشته باشد \cite{rahmawatiagustinaSystematicLiteratureReview2025}. راه‌حل بررسی‌شده، تفکیک مرجع به مرجع ردیابی \LTRfootnote{\lr{Trace Authority (TRA)}} — که تولید شبه‌هویت و ردیابی خودروهای مخرب را بر عهده دارد — و مرکز تولید کلید \LTRfootnote{\lr{Key Generation Center (KGC)}} — که تولید کلیدهای رمزنگاری و پارامترهای امنیتی را انجام می‌دهد — است \cite{rahmawatiagustinaSystematicLiteratureReview2025}. این تفکیک مانع می‌شود یک نهاد واحد بتواند هم هویت کاربران را فاش کند و هم کلیدهای رمزنگاری را در اختیار داشته باشد. نهادهای مستقل دیگری نیز در این مرور درج شده‌اند، از جمله مرجع احراز هویت، مرجع ردیابی قابل اعتماد، مرجع مدیریت اعتماد، مرجع گواهی‌دهی و مرجع ردیابی سازنده \cite{rahmawatiagustinaSystematicLiteratureReview2025}.

مسئلهٔ مرجع حل هویت \LTRfootnote{\lr{Privacy-authority problem}} تنش بین توانایی مرجع در ردیابی هویت واقعی و بی‌نام ماندن خودروهاست و از نبود مرحلهٔ تفکیک هویت در چرخهٔ حیات شبه‌نام ناشی می‌شود \cite{rahmawatiagustinaSystematicLiteratureReview2025}. افزون بر این، در بیشتر سازوکارهای متمرکز واحدهای کنارجاده‌ای به‌عنوان واحدهای کاملاً مورد اعتماد فرض می‌شوند، در حالی که همواره در دسترس نیستند و خودشان آسیب‌پذیرند \cite{farsimadanReviewSecurityChallenges2025}. این حملات زیرساختی و اعتمادی، انگیزهٔ اصلی روی آوردن به سازوکارهای غیرمتمرکز و مدیریت اعتماد توزیع‌شده‌ای است که در فصل‌های ۷ و ۸ بررسی می‌شود.
